\subsection{More native witnesses do not imply faster recognition}

The isolated benchmark runs complete recognition with the native branch
forced, earlier shortcuts disabled, a twelve-second allowance and a
20,000-object fallback allowance. Each timed call includes an additional
independent replay and hash of any returned native proof. An inconclusive
native attempt continues into reduced Khovanov recognition. Shellings-off
and shellings-on each have an A/A control; the four arms run in randomized
order after one warm-up round per case. No tests, audits, builds or commits
overlap either timing run. All per-arm result signatures, work counts and
proof hashes are stable within each run.

The first five-round experiment has 200 measured calls and 40 warm-ups and
takes 13.736 seconds. A/A ratios near 1.09 on one case justify a nine-round
repeat, which has 360 measured calls and 40 warm-ups and takes 25.931 seconds.
All 560 measured calls and 80 warm-ups complete with the expected verdict.
The final primary-arm medians and median paired ratios are:
\begin{center}
\small
\begin{tabular}{lrrrrr}
\toprule
Source & Off ms & On ms & Off/on & Off A/A & On A/A \\
\midrule
circle-9 & 68.402 & 113.360 & 0.598 & 1.046 & 1.021 \\
random-29 & 32.534 & 58.229 & 0.566 & 1.000 & 0.998 \\
optimized-positive & 73.749 & 72.677 & 1.028 & 1.009 & 1.012 \\
random-142 & 111.458 & 140.961 & 0.800 & 1.008 & 0.997 \\
random-1 & 30.707 & 43.739 & 0.703 & 0.989 & 0.991 \\
random-25 & 24.638 & 46.146 & 0.607 & 1.015 & 1.005 \\
random-81 & 76.981 & 129.804 & 0.593 & 0.998 & 1.000 \\
random-176 & 46.083 & 76.112 & 0.606 & 1.007 & 0.997 \\
trefoil & 20.296 & 24.763 & 0.796 & 1.003 & 0.983 \\
genus-one-miss & 22.174 & 38.669 & 0.567 & 0.997 & 1.003 \\
\bottomrule
\end{tabular}

\end{center}
A ratio greater than one favors shellings. Nine of the ten selected
pipelines are slower with shellings, with paired off/on ratios from 0.566
to 0.800. The original optimized-disc fixture is near parity at 1.028;
its lower native guard count does not translate into a substantial total
speedup because the new source proof also authenticates and replays the
simplification. Final A/A medians range from 0.983 to 1.046. Absolute timings
also drift across cases: the identical diagrams labelled 25 and 176 are timed
separately and have different millisecond medians, but similar paired ratios.
These measurements support the default decision, not a universal timing law.

All four added corpus positives switch from the reduced-Khovanov fallback to
a native annulus proof, yet take longer in this measured full-verification
pipeline. The additional three-crossing unknot given by the two-strand braid
$[1,1,-1]$ also gets a native annulus witness after six shellings: tree trial two
produces a connected 58-piece surface with two boundary circles. Its full
proof is retained separately, and Regina confirms the surface summary. The
new total pipeline is slower. Earlier negative results for raw and minimum-span surfaces on
that example therefore remain restricted-family statements. The trefoil
continues through the fallback and is correctly recognized as knotted.
The easy circle and raw-annulus controls incur source-replay overhead on a
path that already succeeded before simplification.

Shellings remain an explicit research and proof-construction option. Their
independently replayable triangulation changes broaden the native family and
can improve bounded-search completion, but the measured corpus does not
justify enabling them by default. Future work can examine proof-checking
costs and progress-based selection between triangulations while retaining
independent source authentication. Neither a short shelling trace nor a
polynomial shelling algorithm bounds how many triangulations or normal
surface families general unknot recognition would need to explore.
The general subexponential and quasipolynomial goals remain open here.

The reproduction driver is \code{normal\_orbit\_research.shellings}.
Select \code{audit} mode, or run \code{recognize --rounds 9} for timings. The implementation is
revision \code{1eb57f251}, published after merging concurrent main-branch
research in \code{77347b0b5}. The incoming unknot inventory and
\code{openai/math} head \code{fd4aeeb2e} were rechecked and are unchanged.
